\subsection{The five leading-profile pressure interactions}
\label{subsec:leading-profile-pressure}

The remaining linear interactions in equation (28) appear at first to
involve an unbounded radial derivative of $V_k$. Their ordered sum has
a cancellation which is more useful than estimating those derivatives
separately. The reduction is exact and uses no assumptions about the
size or sign of the leading profiles.

Put $v=V_0/X$, let $F=U_k\in I_k$, and write $V_k=XW_k$.
Equation (24) gives, with $B=2\eta/L$, $G=d/L$,
\begin{equation}\label{eq:leading-W}
 W_k=BF-B(D+\lambda_k)A_XF-G\partial_\eta A_XF.
\end{equation}
In particular $W_k\in I_k$ coefficientwise whenever $F\in I_k$;
this is a valuation statement and not a claim that the unregularised
map $F\mapsto W_k$ is bounded in the analytic norm. Define
\begin{equation}\label{eq:leading-K0}
 K_0=G\partial_\eta v-B(v+D_Xv).
\end{equation}
The exact chain rule (18) yields
\[
 \frac{Z_{0,k}V_k}{X}
  =G\partial_\eta W_k+B(\lambda_k-1-D_X)W_k,
 \qquad \frac{Z_{0,0}V_0}{X}=K_0.
\]

\begin{proposition}[Radial cancellation and exact pressure primitive]
\label{prop:leading-cancellation}
For every $k\ge1$, the sum of the five leading-profile terms in
$\Omega_k$, namely
\[
 S_k=V_0V_k'+V_kV_0'-\frac{V_0V_k}{X}
     +U_0Z_{0,k}V_k+U_kZ_{0,0}V_0,
\]
admits the exact representation
\begin{equation}\label{eq:leading-source}
 \frac{S_k}{X}=(D_X+1)(vW_k)
     +U_0\bigl[G\partial_\eta+B(\lambda_k-1-D_X)\bigr]W_k
     +FK_0.
\end{equation}
Consequently its contribution to the order-$(k+1)$ pressure is
\begin{equation}\label{eq:leading-primitive}
 \Pi_{k+1}^{\mathrm{lead}}
 =-\frac12 XvW_k-\frac12\int_0^X
 \left\{U_0\bigl[G\partial_\eta+B(\lambda_k-1-D_X)\bigr]W_k
     +FK_0\right\}(x,\eta)\,dx.
\end{equation}
After the shifted axial operator and the axis-regular inverse, its
contribution to the next axial coefficient is precisely
\begin{equation}\label{eq:leading-pressure-operator}
 \mathcal P_{k+1}^{\mathrm{lead}}F
 =-\frac14 J_1 Z_{-2A,k+1}
 \left[XvW_k+\int_0^X
 \left\{U_0\bigl[G\partial_\eta+B(\lambda_k-1-D_X)\bigr]W_k
   +FK_0\right\}dx\right].
\end{equation}
If $v,U_0,K_0$ are regular at the axis, then
$\mathcal P_{k+1}^{\mathrm{lead}}(I_k)\subseteq I_{k+2}$
coefficientwise.
\end{proposition}
\begin{proof}
Since $V_0=Xv$ and $V_k=XW_k$, direct differentiation gives
\[
 V_0V_k'+V_kV_0'-\frac{V_0V_k}{X}
   =X(D_X+1)(vW_k)=X\partial_X(XvW_k).
\]
This proves the first term of \eqref{eq:leading-source}. The other
two terms follow from the displayed $Z$ identities. Equation (27)
contributes $-S_k/(2X)$ to $\Pi_{k+1}'$.
Integrating from the axis, where $XvW_k$ vanishes, gives
\eqref{eq:leading-primitive}. Equation (26), with its factor
$\frac12J_1 Z_{-2A,k+1}$, gives
\eqref{eq:leading-pressure-operator}.
Finally $W_k\in I_k$, and multiplication by axis-regular fixed
profiles preserves radial valuation. The bracket in
\eqref{eq:leading-pressure-operator} lies in $I_{k+1}$;
$J_1$ raises degree by one, proving the inclusion.
\end{proof}

The cancellation has a second useful consequence: the three source
families have different \emph{input} derivative orders after pressure
reconstruction. The $vW_k$ term has at most two derivatives of $F$;
the $U_0Z_{0,k}V_k$ term has at most three; and $FK_0$ has at most one.
Those statements count derivatives of the unknown $F$, not of the
leading profiles. Derivatives of $v,U_0,B,G$ must be included
explicitly when a bounded analytic operator is constructed.

\begin{lemma}[Two-radial-lift scalar derivative budget]
\label{lem:leading-weight-budget}
For $a,b\ge0$, $r=a+1$ and $j\in\{0,1,2,3\}$, the exact
coefficient weights \eqref{eq:weight} satisfy
\begin{equation}\label{eq:leading-mixed-bound}
 \frac{r^{3-j}w_{a,b+j}}{r^2(r+1)^2 w_{a+2,b}}
 \le \frac{54R^2}{r\rho^j}
       \frac{25\,3^{4-j}}{(j+1)^2}.
\end{equation}
\end{lemma}
\begin{proof}
Apply the all-index two-lift inequality
\eqref{eq:alljet} with the auxiliary index $n=a+2$.
It bounds the same weight ratio with the numerator
$(a+3)^{4-j}$ in place of $r^{4-j}$.
Since $r=a+1\le a+3$, division by $r$ gives
\eqref{eq:leading-mixed-bound}.
\end{proof}

\paragraph{What remains to prove.}
The coefficient budget \eqref{eq:leading-mixed-bound} is compatible
with a $1/(k+1)$ estimate for the regularised leading-profile
pressure operator. It is not that operator estimate: each analytic
profile contains a full radial Taylor series, and commuting its
multipliers through the parameter and radial derivatives produces
additional terms. These must be bounded after radial convolution in
the \emph{induced} norm with constants independent of $k$.
In particular, applying a boundedness claim to the separate map
$F\mapsto\partial_\eta A_XF$ would be invalid. The precise remaining
analytic statement is a uniform bound for the full
\eqref{eq:leading-pressure-operator}, with the common holomorphic
radius and multiplication constants of the actual leading profiles
specified. The positive-order quadratic pressure sum and the other
momentum forcing channels are separate subsequent steps; neither the
recurrence (53) nor Theorem~1.1 follows from Proposition
\ref{prop:leading-cancellation} alone.

The accompanying exact symbolic replay compares both sides of the
five-term source identity and its reconstructed pressure coefficient
on 72 independently specified nontrivial cases, and verifies 41,472
rational inequalities of the form \eqref{eq:leading-mixed-bound}.
The corresponding Lean source encodes the finite radial-cancellation,
primitive-factor and valuation lemmas; kernel certification is pending
successful compilation in the pinned GitHub workflow. The complete analytic
operator is not represented as Lean-certified.
